% !TeX root = ../main.tex
\section{Algorytm}

\begin{frame}{Główne fazy}
  \begin{enumerate}[<+->]
    \item Losowe wygenerowanie kandydatów na lidera.
        \begin{enumerate}[<+->]
            \item Jeśli wygenerowany zbiór jest pusty powtarzamy fazę 1.
        \end{enumerate}
    \item Weryfikujemy że został wybrany dokładnie jeden kandydat.
        \begin{enumerate}[<+->]
            \item Któryś z węzłów orientuje się, że jest więcej niż jeden lider.
            \item Wszystkie węzły wiedzą czy wybrano unikatowego lidera
        \end{enumerate}
  \end{enumerate}
  \pause
  Procedura ta jest powtarzana aż zostanie wybrany dokładnie jeden lider.
\end{frame}

\begin{frame}{Faza 1}
    Cała procedura trwa $2N-1 \geq n$ rundy.
    \begin{enumerate}[<+->]
        \item Z prawdopodobieństwem $\frac{c}{N}$ węzeł zostaje kandydatem
        \item Komunikacja:
            \begin{itemize}
                \item Kandydaci inicjują komunikację wysyłając bit na wyjście
                \item Pozostałe węzły tylko przekazują wiadomości dalej
            \end{itemize}
        \item Węzeł wie, że istnieje jakiś kandydad dokładnie gdy zobaczył jakąś wiadomość
    \end{enumerate}
\end{frame}

\begin{frame}{Faza 2.1}
    \textbf{Cel: Jeśli jest więcej niż jeden kandydat to któryś węzeł musi to zauważyć}

    \begin{columns}[T,onlytextwidth]
    \begin{column}{0.6\textwidth}
    \uncover<2->{
    Przez dokładnie $N-1$ rund węzły komunikują się w następujący sposób:
    \begin{itemize}
        \item Kandydaci - inicjują komunikację wysyłając bit. Nie przekazują dalej odczytanych wiadomości.
        \item Pozostałe węzły - przekazują dalej wiadomości
    \end{itemize}
    }
    \uncover<3->{
    Jeśli kandydat otrzymał bit to wie że jest więcej niż jeden kandydat (bo $n > N-1$).
    }

    \uncover<4->{
    Jeśli jest więcej niż jeden kandydat to muszą być dwaj kandydaci w odległości co najwyżej $N-1$.
    }
    \end{column}

    \begin{column}{0.36\textwidth}
      \centering
      % N = 3, n = 5; candidates v_1, v_3 at distance 2 = N - 1.
      % v_1's bit reaches v_3 after 2 rounds; v_3's bit stops at v_5.
      \begin{tikzpicture}
        \foreach \i in {1,...,5} {
          \node[proc] (v\i) at ({90 - 72*(\i-1)}:1.5cm) {$v_\i$};
        }
        \foreach \i/\j in {1/2, 2/3, 3/4, 4/5, 5/1} {
          \draw[dlink] (v\i) to[bend left=20] (v\j);
        }
        % candidates
        \node[proc, active] at (v1) {$v_1$};
        \node[proc, active] at (v3) {$v_3$};
        % messages, labelled with the round in which they are sent
        \draw<2->[msg] (v1) to[bend left=20] node[auto, font=\scriptsize] {1} (v2);
        \draw<2->[msg] (v2) to[bend left=20] node[auto, font=\scriptsize] {2} (v3);
        \draw<2->[msg] (v3) to[bend left=20] node[auto, font=\scriptsize] {1} (v4);
        \draw<2->[msg] (v4) to[bend left=20] node[auto, font=\scriptsize] {2} (v5);
        % the candidate that detects another candidate
        \node<3->[proc, aware] at (v3) {$v_3$};
      \end{tikzpicture}

      \medskip
      {\scriptsize $N = 3$, $n = 5$\\ liczby przy strzałkach: numer rundy\par}
      \smallskip
      {\scriptsize
        \tikz\node[proc, active, minimum size=3mm]{}; kandydat\quad
        \only<3->{\tikz\node[proc, aware, minimum size=3mm]{}; kandydat, który wie}
      \par}
    \end{column}
    \end{columns}
\end{frame}

\begin{frame}{Faza 2.2}
    \textbf{Cel: wszyscy wiedzą czy się udało} 

    \pause
    Przez dokładnie $2N-1$ rund węzły komunikują się w następujący sposób:
    \begin{itemize}
        \item Kandydaci którzy wiedzą że jest więcej niż jeden kandydat - inicjują komunikację wysyłając bit. Nie przekazują dalej odczytanych wiadomości.
        \item Pozostałe węzły - przekazują dalej wiadomości
    \end{itemize}
    \pause 
    Procedura zakończyła się sukcesem dokładnie gdy zaobserwowaliśmy jakąś wiadomość
\end{frame}